\SGAchapter{XI}{例子与补充}{Exemples et compléments}
\SourcePage{235}
\SGAsection{1}{射影空间与单有理簇}
\begin{SGAresult}{命题}{1.1}
设 $k$ 为\Term{algebraically-closed-field}{代数闭域}，$X=\mathbf P_k^r$ 为 $k$ 上的 $r$ 维\Term{projective-space}{射影空间}。则 $X$ \Term{simply-connected}{单连通}，即 $\pi_1(X)=0$。
\end{SGAresult}
$r=0$ 时显然。若 $r=1$，需要证明：$X=\mathbf P_k^1$ 的任意非空\Term{connected}{连通}\Term{etale-covering}{平展覆盖} $X'$ 都满足 $X'\xrightarrow{\sim}X$。若 $g,g'$ 分别为 $X,X'$ 的\Term{genus}{亏格}，则\Term{genus}{亏格}公式给出
\[
1-g'=d(1-g),
\]
其中 $d$ 为 $X'$ 在 $X$ 上的次数。由于 $g=0$，故 $1-g'=d$；而 $g'\geq0$，所以必有 $d=1$，从而 $X'\xrightarrow{\sim}X$。$r\geq2$ 时对 $r$ 归纳，假设对 $r'<r$，$\mathbf P^{r'}$ 均\Term{simply-connected}{单连通}。将此用于 $\mathbf P^r$ 的一个超平面，再用（X 2.10），即得 $\mathbf P^r$ \Term{simply-connected}{单连通}。另一个证明是：由（X 1.7），
\[
\pi_1(\mathbf P^1\times\cdots\times\mathbf P^1)
=\pi_1(\mathbf P^1)\times\cdots\times\pi_1(\mathbf P^1),
\]
所以 $\mathbf P^1$ 的\Term{simply-connected}{单连通}性蕴含 $(\mathbf P^1)^r$ \Term{simply-connected}{单连通}，再由\Term{birational-invariance}{基本群的双有理不变性}（X 3.4），得到 $\mathbf P^r$ \Term{simply-connected}{单连通}。这个证明还给出更一般的结论：

\begin{SGAresult}{推论}{1.2}
设 $X$ 为\Term{algebraically-closed-field}{代数闭域} $k$ 上\Term{proper}{固有}且\Term{normal}{正规}的\Term{scheme}{概形}。若 $X$ 为\Term{rational-variety}{有理簇}，即 $X$ 为\Term{integral}{整}的，且其\Term{function-field}{函数域}是 $k$ 的\Term{purely-transcendental}{纯超越扩张}，则 $X$ \Term{simply-connected}{单连通}。
\end{SGAresult}
这一结果特别适用于 Grassmann 簇；更一般地，适用于形如 $G/H$ 的簇，其中 $G$ 为 $k$ 上的\Term{connected}{连通}线性群，而 $H$ 为包含 $G$ 的一个 Borel 子群的代数子群。

回忆：$k$ 上的\Term{unirational-variety}{单有理簇}，指 $k$ 上的一个\Term{proper}{固有}整\Term{scheme}{概形}，其\Term{function-field}{函数域} $K$ 包含于 $k$ 的某个\Term{purely-transcendental}{纯超越扩张} $K'$，且 $K'/K$ \Term{finite}{有限}，即 $K'$ 与 $K$ 在 $k$ 上的\Term{transcendence-degree}{超越次数}相同。等价地，存在\Term{dominant-rational-map}{支配有理映射} $f:\mathbf P_k^r\to X$，其中 $r=\dim X$。若 $X$ \Term{normal}{正规}，由 X 3.4 之前的讨论可知：对 $X$ 的任意\Term{connected}{连通}\Term{etale-covering}{平展覆盖} $X'$，若相应的域\Term{extension}{扩张}为 $L/K$，则 $K'$ 上的代数 $L\otimes_KK'$ 在模型 $\mathbf P^r$ 上\Term{unramified}{非分歧}。
\SourcePage{236}
于是由 1.1，该代数\Term{completely-split}{完全分裂}；这表明 $L$ 与 $K'/K$ 的某个子\Term{extension}{扩张}在 $K$ 上\Term{isomorphism}{同构}。结合 V 8.2，得到：

\begin{SGAresult}{推论}{1.3}
\Term{algebraically-closed-field}{代数闭域}上的\Term{normal}{正规}\Term{unirational-variety}{单有理簇}的\Term{fundamental-group}{基本群}是\Term{finite-group}{有限群}。
\end{SGAresult}
（注意，在“单有理”的定义中，其实无须要求 $K'/K$ \Term{finite}{有限}。）

\textbf{注记 1.4。} 本节结果当然都是熟知的。另一方面，J.-P. Serre 已证明 [10]：若 $X$ 为特征零\Term{algebraically-closed-field}{代数闭域}上的\Term{smooth}{光滑}\Term{projective}{射影}\Term{unirational-variety}{单有理簇}，则 $X$ \Term{simply-connected}{单连通}。他的证明使用 Hodge 对称性定理，因而是一个超越证明。

〔2003 年增补（MR）。限于特征 $p\geq0$ 的\Term{algebraically-closed-field}{代数闭域} $k$。在特征 $p>0$ 时，上述 $k$-\Term{unirational-variety}{单有理簇}的定义，是\Term{weakly-unirational}{弱单有理} $k$-簇的定义；与之相对的是\Term{strongly-unirational}{强单有理} $k$-簇的定义，后者还要求 $K'/K$ 为\Term{separable-extension}{可分扩张}。

在维数 2 的情形，存在\Term{smooth}{光滑}\Term{projective}{射影}\Term{weakly-unirational}{弱单有理}曲面，其\Term{fundamental-group}{基本群}非\Term{trivial-bundle}{平凡}（由 1.3，该群\Term{finite}{有限}），因而这些曲面并不有理。参见 T. Shioda，\textit{On unirationality of supersingular surfaces}，Math. Ann. 225（1977），155—159 页。反之，\Term{strongly-unirational}{强单有理}曲面必有理：这是 Castelnuovo 有理性判据的推论，该判据由 O. Zariski 推广到特征 $p>0$；参见 J.-P. Serre，Bourbaki 讨论班第 146 讲（1957），及其《全集》第 1 卷，第 467 页。

在复数域 $\mathbf C$ 上，已知维数 $\geq3$ 的\Term{smooth}{光滑}\Term{projective}{射影}簇中有单有理但非有理的例子。参见 P. Deligne，\textit{Variétés unirationnelles non rationnelles [d'après M. Artin et D. Mumford]}，Bourbaki 讨论班第 402 讲，1971—1972 年，Lecture Notes 第 317 卷。特别地，\Term{projective-space}{射影空间} $\mathbf P^4$ 中的\Term{smooth}{光滑}三次超曲面就是这样的例子。其中一些例子可推广到正特征，给出\Term{strongly-unirational}{强单有理}但非有理的簇。参见 J. P. Murre，\textit{Reduction of the proof of the non-rationality of a non singular cubic threefold to a result of Mumford}，Compositio Math. 27（1973），63—82 页。

设 $V$ 为\Term{normal}{正规}、整的\Term{projective}{射影} $k$-簇。若存在\Term{finite-type}{有限型}整 $k$-\Term{scheme}{概形} $T$ 及 $k$-\Term{morphism}{态射} $F:T\times\mathbf P^1\to V$，使\Term{morphism}{态射}
\[
\begin{aligned}
T\times\mathbf P^1\times\mathbf P^1&\longrightarrow V\times V,\\
(t,u,u')&\longmapsto\bigl(F(t,u),F(t,u')\bigr)
\end{aligned}\tag{1}
\]
为\Term{dominant}{支配}的，则称 $V$ 为\Term{rationally-connected}{有理连通}的。
\SourcePage{237}
特别地，$V$ 的两个充分一般的有理点可由一条有理曲线连结。这一术语的依据是：若 $V$ \Term{rationally-connected}{有理连通}，则任意两个有理点都可由\Term{finite}{有限}条有理曲线组成的链连结。当 $k$ 的特征为 $p>0$ 时，若进一步要求映射（1）\Term{generically-smooth}{一般光滑}，便得到\Term{separably-rationally-connected}{可分有理连通}簇的概念。因此，\Term{unirational-variety}{单有理簇}\Term{rationally-connected}{有理连通}；在特征 $p>0$ 时，\Term{strongly-unirational}{强单有理}簇\Term{separably-rationally-connected}{可分有理连通}。J. Kollár 证明了\Term{separably-rationally-connected}{可分有理连通}簇的代数\Term{fundamental-group}{基本群}\Term{trivial-bundle}{平凡}。参见 O. Debarre，\textit{Variétés rationnellement connexes [d'après T. Graber, J. Harris, J. Starr et A. J. de Jong]}，Bourbaki 讨论班第 906 讲，2001—2002 年。〕

\SGAsection{2}{阿贝尔簇}
设 $k$ 为\Term{algebraically-closed-field}{代数闭域}，$A$ 为 $k$ 上的\Term{abelian-variety}{阿贝尔簇}，即 $k$ 上\Term{proper}{固有}、\Term{smooth}{光滑}且\Term{connected}{连通}的\Term{group-scheme}{群概形}；再设 $G$ 为 $k$ 上\Term{finite-type}{有限型}的交换\Term{group-scheme}{群概形}。以 $\operatorname{Ext}(A,G)$ 表示 $A$ 被 $G$ 的交换\Term{extension}{扩张}的等价类所成的群，以 $H^1(A,G)$ 表示 $A$ 上以 $G$ 为群的\Term{principal-bundle}{主丛}的等价类所成的群（参见下面第 4 节），并考虑典范\Term{homomorphism}{同态}
\[
\operatorname{Ext}(A,G)\longrightarrow H^1(A,G).
\]
Serre 的一个论证 [5，第 VII 章，定理 5] 表明，该\Term{homomorphism}{同态}为\Term{injectivity}{单射}，其像正是 $H^1(A,G)$ 的“\Term{primitive-element}{本原元}”全体，即满足
\[
\pi^*(\xi)=\operatorname{pr}_1^*(\xi)+\operatorname{pr}_2^*(\xi)
\]
的元素 $\xi$；这里 $\operatorname{pr}_i$ 为 $A\times A$ 到 $A$ 的两个投影，$\pi:A\times A\to A$ 为 $A$ 的合成律。（Serre 的定理只对线性、\Term{connected}{连通}且当然在 $k$ 上\Term{smooth}{光滑}的 $G$ 陈述。不过，简化其论证的前半部分便可看出，这些限制并无必要：只须注意，由 $A$ 到 $k$ 上\Term{finite-type}{有限型}\Term{group-scheme}{群概形} $E$ 的任意保单位元的\Term{morphism}{态射}，必为\Term{group-homomorphism}{群同态}；再将此用于 $A$ 被 $G$ 的\Term{extension}{扩张} $E$ 在 $A$ 上的\Term{section}{截面}即可。）

将此结果用于 $G$ 为 $k$ 上\Term{finite}{有限}\Term{separable}{可分}群的情形，即 $G$ 为一个通常的\Term{finite-group}{有限群}，并假定其交换。利用 $\pi_1(A\times A)\xrightarrow{\sim}\pi_1(A)\times\pi_1(A)$（X 1.7），并对任意代数\Term{scheme}{概形} $X$，特别对 $X=A$ 或 $A\times A$，将 $H^1(X,G)$ 解释为 $\operatorname{Hom}(\pi_1(X),G)$，可知 $H^1(A,G)$ 的每个类均本原。因此有\Term{isomorphism}{同构}
\[
\operatorname{Ext}(A,G)\xrightarrow{\sim}H^1(A,G).
\]
\SourcePage{238}
换言之，$A$ 的任意以交换群 $G$ 为结构群、并在 $A$ 的原点上方带有标记点的\Term{principal-covering}{主覆盖}，都唯一地带有一个以该标记点为原点的代数群结构，使 $A'\to A$ 为代数\Term{group-homomorphism}{群同态}。特别地，若 $A'$ \Term{connected}{连通}，则它也是一个\Term{abelian-variety}{阿贝尔簇}，并与 $A$ \Term{isogenous}{同源}。

另一方面，带点代数\Term{scheme}{概形} $X$ 到群的\Term{functor}{函子} $X\mapsto\pi_1(X)$ 与乘积交换（IX 1.7），因此将前一\Term{category}{范畴}中的群变为群\Term{category}{范畴}中的群，即交换群。所以，若 $A$ 为\Term{abelian-variety}{阿贝尔簇}，则 $\pi_1(A)$ 是交换群。因此，要确定 $\pi_1(A)$，只须确定\Term{finite}{有限}交换群 $G$ 上的\Term{functor}{函子}
\[
G\longmapsto H^1(A,G)=\operatorname{Hom}(\pi_1(A),G).
\]
最后回忆，对每个整数 $n>0$，$A$ 上的乘 $n$ \Term{homomorphism}{同态}
\[
A\xrightarrow{n}A
\]
均为\Term{surjectivity}{满射}，因而其核\Term{finite}{有限}，即该\Term{homomorphism}{同态}为\Term{isogeny}{同源}。由此可得，任意\Term{isogenous}{同源} $A'\to A$ 都是这种\Term{isogenous}{同源}的商。结合这些事实与标准论证（例如见 [6]），得到：

\begin{SGAresult}{定理（Serre–Lang）}{2.1}
设 $A$ 为\Term{algebraically-closed-field}{代数闭域} $k$ 上的\Term{abelian-variety}{阿贝尔簇}。对每个整数 $n>0$，设 $K_n$ 为 $A$ 上乘 $n$ \Term{homomorphism}{同态}的核 ${}_nA$ 所对应的通常\Term{finite-group}{有限群}；再对每个素数 $\ell$ 置
\[
T_\ell(A)=\varprojlim_r K_{\ell^r},
\]
并置
\[
T_{\boldsymbol\cdot}(A)=\prod_\ell T_\ell(A)=\varprojlim_n K_n.
\]
这里，当 $m$ 为 $n$ 的倍数，$m=ns$ 时，以乘 $s$ \Term{homomorphism}{同态}将 $K_m$ 映到 $K_n$。则 $\pi_1(A)$ 与 $T_{\boldsymbol\cdot}(A)$ 典范\Term{isomorphism}{同构}；因此，对每个素数 $\ell$，$\pi_1(A)$ 的 $\ell$-准素分量与 $T_\ell(A)$ 典范\Term{isomorphism}{同构}。
\end{SGAresult}
注意，这些\Term{isomorphism}{同构}对变化的 $A$ 具有\Term{functorial}{函子性}。模 $T_\ell(A)$ 称为\Term{abelian-variety}{阿贝尔簇} $A$ 的 $\ell$ 进\Term{tate-module}{Tate 模}。它是 $A$ 的\Term{additive-functor}{加性函子}，特别地，它给出 $A$ 的自\Term{homomorphism}{同态}环 $\operatorname{Hom}(A,A)$ 在 $T_\ell(A)$ 上的一个表示，称为\Term{weil-adic-representation}{Weil 的 $\ell$ 进表示}，在\Term{abelian-variety}{阿贝尔簇}理论中具有重要作用，例如见 [4，第 VII 章]。定理 2.1 将它解释为 $A$ 的 $\ell$ 进同调群
\[
H_1(A,\mathbf Z_\ell)=\pi_1(A)_\ell
\]
上的自然表示；尤其从 Lefschetz 公式的观点看 [4，第 V 章]，这种解释显然在概念上更为令人满意。这里回忆 Weil 关于 $T_\ell(A)$ 结构的结果：
\SourcePage{239}
\begin{enumerate}[label=\alph*)]
\item 若 $n$ 与 $\operatorname{car}(k)$ 互素，则 $K_n$ 为 $\mathbf Z/n\mathbf Z$ 上秩 $2\dim A$ 的\Term{free-module}{自由模}。因此，若素数 $\ell\ne\operatorname{car}(k)$，则 $T_\ell(A)$ 为 $\ell$ 进整数环 $\mathbf Z_\ell$ 上秩 $2\dim A$ 的\Term{free-module}{自由模}。
\item 若 $n$ 为 $\operatorname{car}(k)=p$ 的幂，则 $K_n$ 为 $\mathbf Z/n\mathbf Z$ 上秩 $\nu\leq\dim A$ 的\Term{free-module}{自由模}，且 $\nu$ 与 $n$ 无关。因此 $T_p(A)$ 为 $p$ 进整数环 $\mathbf Z_p$ 上秩 $\nu\leq\dim A$ 的\Term{free-module}{自由模}。
\end{enumerate}
这表明，在这里发展的\Term{fundamental-group}{基本群}理论中，变化的\Term{abelian-variety}{阿贝尔簇}的\Term{fundamental-group}{基本群}并不随参数规则地变化：当参数 $t$ 对应的\Term{residue-characteristic}{剩余特征}为 $p$ 时，它的 $p$-准素分量可能突然缩小。这一现象最熟知的例子出现在椭圆曲线中。不过应注意，无论 $n$ 是否与特征互素，$A$ 上乘 $n$ \Term{homomorphism}{同态}的真正的核 ${}_nA$ 都是 $k$ 上次数 $n^{2g}$ 的\Term{finite}{有限}\Term{group-scheme}{群概形}，其中 $g=\dim A$；若 $n$ 是 $p=\operatorname{car}(k)$ 的倍数，该\Term{group-scheme}{群概形}在 $k$ 上不\Term{separable}{可分}。而且，当 $A$ 在\Term{abelian-variety}{阿贝尔簇}族中变化，即给定基\Term{scheme}{概形} $S$ 上的阿贝尔\Term{scheme}{概形} $A$ 时，更一般地可证 ${}_nA$ 为 $S$ 上次数 $n^{2g}$ 的\Term{finite}{有限}\Term{flat}{平坦}\Term{group-scheme}{群概形}。也就是说，只要计入核 ${}_nA$ 的\Term{infinitesimal-part}{无穷小部分}，那么无论 $n$ 为何，它们都规则地变化。这提示我们，\Term{abelian-variety}{阿贝尔簇} $A$ 的“真正的”\Term{fundamental-group}{基本群}应当是投代数群
\[
\varprojlim_n {}_nA,
\]
即 $k$ 上\Term{finite-group}{有限群}的形式\Term{projective-limit}{投极限}。所谓代数\Term{scheme}{概形} $X$ 的“真正的\Term{fundamental-group}{基本群}”，应理解为：分类 $X$ 的所有\Term{principal-covering}{主覆盖}的投群，其中结构群允许为 $k$ 上任意\Term{finite-group}{有限群} $G$，不必在 $k$ 上\Term{separable}{可分}。例如，这样就可以借助 $\operatorname{Hom}(A,A)$ 在 $A$ 的真正\Term{fundamental-group}{基本群}的 $p$-准素分量上的表示，恢复 Weil 用 $\ell\ne p$ 定义的特征多项式，而这比 Serre 的构造 [8] 更自然。

\SGAsection{3}{投影锥与 Zariski 的例子}
为简便，仍设 $k$ 为\Term{algebraically-closed-field}{代数闭域}，$V$ 为\Term{connected}{连通}的\Term{projective}{射影} $k$-\Term{scheme}{概形}，它是 $\mathbf P_k^r$ 的一个\Term{closed-subscheme}{闭子概形}；愿意的话，也可假设它\Term{nonsingular}{非奇异}。设 $Y=\widehat C$ 为 $V$ 的\Term{projective}{射影}\Term{projecting-cone}{投影锥}，$y_0$ 为其顶点；$X=\widehat C_V$ 为与 $O_V(1)$ 关联的\Term{vector-bundle}{向量丛} $C_V=\mathbf V(O_V(1))$ 的通常\Term{projective-closure}{射影闭包}。再设
\[
f:X\longrightarrow Y
\]
为典范\Term{morphism}{态射}，它将 $X$ 中 $C_V$ 的\Term{zero-section}{零截面} $X_0$ 收缩为一点（EGA II 8.6.4）。由于 $X$ 为 $V$ 上以 $\mathbf P^1$ 为\Term{fiber}{纤维}的\Term{locally-trivial-bundle}{局部平凡丛}，而这些\Term{fiber}{纤维}\Term{simply-connected}{单连通}，由（X 1.4），\Term{morphism}{态射} $p:X\to V$ 诱导\Term{isomorphism}{同构}
\SourcePage{240}
\[
\pi_1(X)\xrightarrow{\sim}\pi_1(V).
\]
由于 $p$ 诱导\Term{isomorphism}{同构} $X_0\to V$，可知 $X$ 的一个\Term{etale-covering}{平展覆盖}\Term{completely-split}{完全分裂}，当且仅当其在 $X_0$ 上的限制\Term{completely-split}{完全分裂}。对 $Y$ 的任意\Term{etale-covering}{平展覆盖} $Y'$，\Term{inverse-image}{逆像} $X'=X\times_Y Y'$ 都是 $X$ 的\Term{etale-covering}{平展覆盖}，且在\Term{fiber}{纤维} $X_0$ 上\Term{completely-split}{完全分裂}，因而\Term{trivial-bundle}{平凡}。又因\Term{homomorphism}{同态} $\pi_1(X)\to\pi_1(Y)$ 为\Term{surjectivity}{满射}（IX 3.4），得到
\[
\pi_1(Y)=(e).
\]
换言之，每个\Term{projective}{射影}\Term{projecting-cone}{投影锥}都\Term{simply-connected}{单连通}。（注意，在特征 0 时，取 $Y$ 为\Term{affine}{仿射}\Term{projecting-cone}{投影锥}，同样的结果也成立。）

现在假设 $V$ \Term{regular}{正则}，即在 $k$ 上\Term{smooth}{光滑}。则 $X$ 也\Term{regular}{正则}，并且对 $V$ 选择适当的\Term{projective}{射影}嵌入，可使\Term{projecting-cone}{投影锥} $Y$ \Term{normal}{正规}。若 $V$ 不\Term{simply-connected}{单连通}，因而 $X$ 不\Term{simply-connected}{单连通}，取 $X$ 的一个非\Term{trivial-bundle}{平凡}\Term{connected}{连通}\Term{etale-covering}{平展覆盖} $X'$。由于 $X$ 在 $y\in Y$、$y\ne y_0$ 处的\Term{fiber}{纤维}都是单点，$X'$ 在这些\Term{fiber}{纤维}上、特别是在\Term{generic-fiber}{一般纤维}上的限制\Term{trivial-bundle}{平凡}；然而，$X$ 并非来自 $Y$ 的某个\Term{etale-covering}{平展覆盖}的\Term{inverse-image}{逆像}，因为 $Y$ \Term{simply-connected}{单连通}，否则 $X'$ 将\Term{completely-split}{完全分裂}。这表明，若将（X 1.3、1.4）中 $f$ \Term{separable}{可分}的假设削弱为其\Term{fiber}{纤维}是 $k(s)$ 上\Term{separable}{可分}、甚至\Term{smooth}{光滑}的代数\Term{scheme}{概形}，结论就不再成立。同样，对 $y\ne y_0$，\Term{geometric-fiber}{几何纤维} $X_{\bar y}$ 的\Term{fundamental-group}{基本群}显然为 $(e)$，因为这些\Term{fiber}{纤维}是单点；而 $\pi_1(X_0)\ne e$，所以\Term{semicontinuity}{半连续性}定理（X 2.4）对 $f$ 也不成立。

最后给出 Zariski 指出的例子：将上述定理中 $f$ \Term{separable}{可分}的假设替换为 $f$ \Term{flat}{平坦}时，这些定理仍会失效。设 $f:X\to Y$ 为\Term{nonsingular}{非奇异}\Term{projective}{射影}曲面到有理直线 $Y=\mathbf P^1$ 的\Term{morphism}{态射}，满足 $K=k(X)$ 是 $k(f)$ 的“\Term{regular}{正则}”\Term{extension}{扩张}，即\Term{primary-extension}{准素扩张}且\Term{separable}{可分}（等价地，几何\Term{generic-fiber}{一般纤维}\Term{connected}{连通}且\Term{separable}{可分}），并且\Term{divisor}{除子} $(f)=X_0-X_\infty$ 为某个\Term{divisor}{除子}的 $n$ 倍，其中整数 $n$ 与特征互素。这样的例子可在任意特征下构造。设 $X'$ 为 $X$ 在 $K(f^{1/n})$ 中的\Term{normalization}{正规化}，其中 $K=k(X)$ 为 $X$ 的\Term{function-field}{函数域}。关于 $(f)$ 的假设蕴含 $X'$ 在 $X$ 上\Term{etale}{平展}。再设 $Y'$ 为 $Y$ 在 $k(t)(t^{1/n})$ 中的\Term{normalization}{正规化}；它在 $Y$ 上恰好于 $t=0,\infty$ 处\Term{ramified}{分歧}，而 $X'|_{f^{-1}(U)}$ \Term{isomorphism}{同构}于 $Y'|_U$ 的\Term{inverse-image}{逆像}。特别地，$X'$ 在 $X/Y$ 的几何\Term{generic-fiber}{一般纤维}上的限制\Term{completely-split}{完全分裂}。
\SourcePage{241}
但是 $X'$ 不\Term{isomorphism}{同构}于 $Y$ 的任何\Term{etale-covering}{平展覆盖}的\Term{inverse-image}{逆像}，因为该\Term{covering}{覆盖}显然必为 $Y'$，而 $Y'$ 在 $Y$ 上\Term{ramified}{分歧}，矛盾。\footnote{从“\Term{etale-topology}{平展拓扑}”的观点（SGA 4 VII）看，可以注意到在此例中 $R^1(f_{\text{ét}})_*(\mathbf Z/n\mathbf Z)$ 在 $S$ 上“非\Term{separated}{分离}”。}

下面给出一种简单的实现方法（归于 Serre），其思想来自 [5，第 20 节]。取与特征不同的素数 $n\geq5$，令 $G=\mathbf Z/n\mathbf Z$ 以四个不同的特征标分别乘四个坐标的方式作用在 $k^4$ 上；\footnote{〔2003 年增补：$k^4$ 表示 $k$ 上的 4 维\Term{affine}{仿射}空间。〕}由于 $n\geq5$，这可以做到。于是 $G$ 作用在\Term{projective-space}{射影空间} $\mathbf P_k^3$ 上，且仅有四个不动点，分别对应于坐标轴。方程
\[
x^n+y^n+z^n+t^n=0
\]
给出的曲面 $X'$ 在 $k$ 上\Term{smooth}{光滑}（由 \Term{jacobian-criterion}{Jacobian 判据}），且不包含这些不动点。由于 $G$ 的阶为素数，故它在 $X$ 上的作用“无不动点”，即 $X$ 是 $X=X'/G$ 的以 $G$ 为群的\Term{principal-covering}{主覆盖}。设 $g=x/y\in k(X')=K'$；若所取特征标为 $\chi^i$（$i=0,1,2,3$），其中 $\chi$ 为本原特征标，则 $g$ 为 $K'/K$ 的 Kummer 生成元，这里 $K=k(X)$。设 $f=g^n\in K$。很容易看出 $K'$ 是 $k(g)$ 的\Term{regular}{正则}\Term{extension}{扩张}，因为以 $U,T,Z$ 为齐次坐标、方程为
\[
T^n+Z^n+(1+g^n)U^n=0
\]
的平面曲线在 $k(g)$ 上\Term{smooth}{光滑}（\Term{jacobian-criterion}{Jacobian 判据}），且任意平面曲线都\Term{connected}{连通}。另一方面，$k(f)=K\cap k(g)$，因为右端是 $k(f)$ 的一个\Term{extension}{扩张}，包含于素次数\Term{extension}{扩张} $k(g)$，并且不同于 $k(g)$，因为 $g\notin K$。这蕴含 $K/k(f)$ 为\Term{regular}{正则}\Term{extension}{扩张}。最后，$X$ 上 $f$ 的\Term{divisor}{除子}为某个\Term{divisor}{除子}的 $n$ 倍：它在 $X'$ 上的\Term{inverse-image}{逆像}为 $g^n$ 的\Term{divisor}{除子}，因而是 $n$ 倍；又因 $X'$ 在 $X$ 上\Term{etale}{平展}，可以将这一性质\Term{descent}{下降}。若\Term{rational-map}{有理映射} $f:X\to\mathbf P^1$ 是一个\Term{morphism}{态射}，也就是说，若 $f$ 的零点\Term{divisor}{除子}和极点\Term{divisor}{除子}不相交，证明就完成了。事实上，仍在 $X'$ 上考察，很容易验证上述两个\Term{divisor}{除子}分别是两条 $k$ 上\Term{smooth}{光滑}曲线的 $n$ 倍，两曲线在一点 $a$ 处横截相交。将 $X$ 换为在 $a$ 处\Term{blow-up}{爆破}所得的\Term{scheme}{概形}，则前述条件——$\operatorname{div}(f)$ 可被 $n$ 整除，且 $k(X_1)=k(X)$ 为 $k(X)$ 的\Term{regular}{正则}\Term{extension}{扩张}——仍然成立，而此时 $f$ 还成为\Term{morphism}{态射} $X_1\to\mathbf P^1$，因而满足所需条件。

\SGAsection{4}{上同调正合列}
设 $S$ 为\Term{prescheme}{预概形}，因而 $X$ 上\Term{prescheme}{预概形}的\Term{category}{范畴} $\mathrm{Sch}_{/S}$ 已确定，其中群的概念也随之确定；我们将这样的群称为 $S$ 上的\Term{group-prescheme}{群预概形}，或简称 $S$-群。为简化叙述并明确讨论对象，以下通常只考虑 $S$ 上\Term{affine}{仿射}且\Term{flat}{平坦}的群，\footnote{事实上，下文只须假设拟\Term{affine}{仿射}，无须\Term{affine}{仿射}，参见原书第 228 页脚注（5）。}这足以满足这里预想的应用。当然，也有许多情形并不满足这两个假设中的任何一个。
\SourcePage{242}
设 $G$ 为这样的 $S$-群，$P$ 为带有 $G$ 右作用的 $S$-\Term{prescheme}{预概形}；特别地，这给出 $\mathrm{Sch}_{/S}$ 中的\Term{morphism}{态射}
\[
\pi:P\times_SG\longrightarrow P,
\]
满足熟知的公理。若以 $\operatorname{pr}_1$ 和 $\pi$ 为两个分量的\Term{morphism}{态射}
\[
P\times_SG\longrightarrow P\times_SP
\]
为\Term{isomorphism}{同构}，则称 $P$ 为 $G$ 下的\Term{formally-principal-homogeneous}{形式主齐性}对象。等价地，对 $\mathrm{Sch}_{/S'}$ 的任意对象 $S'$，集合 $P(S')=\operatorname{Hom}_S(S',P)$ 在群 $G(S')=\operatorname{Hom}_S(S',G)$ 的作用下，或者为空，或者为\Term{principal-homogeneous}{主齐性}的；后者指它\Term{isomorphism}{同构}于以右平移作用于自身的 $G(S')$。若 $P$ \Term{isomorphism}{同构}于以右平移作用于自身的 $G$，则称 $P$ \Term{trivial-bundle}{平凡}；等价地，每个带 $G(S')$ 作用的集合 $P(S')$ 都\Term{trivial-bundle}{平凡}。可以验证，例如仍使用那套转化到集合情形的惯常方法：$P$ \Term{trivial-bundle}{平凡}，当且仅当它\Term{formally-principal-homogeneous}{形式主齐性}并且在 $S$ 上有一个\Term{section}{截面}。后一条件以\Term{category}{范畴}语言表述，就是 $P$ 在 $\mathrm{Sch}_{/S}$ 的\Term{final-object}{终对象} $e=S$ 上有\Term{section}{截面}，即存在 $e\to P$ 的\Term{morphism}{态射}。

为了定义比\Term{formally-principal-homogeneous}{形式主齐性}更强的概念——$G$ 下的\Term{principal-homogeneous-bundle}{主齐性丛} $P$——须先在 $\mathrm{Sch}_{/S}$ 中指定一族用于\Term{descent}{下降}的\Term{morphism}{态射}。它们将充当“\Term{localization}{局部化}\Term{morphism}{态射}”，用来使丛“\Term{trivial-bundle}{平凡}化”。最合适的选择随情境而异，没有一种选择包含其余所有选择。\footnote{参见 SGA 3 IV，特别是第 4 节。}在这里，采用下述定义较为方便。

\begin{SGAresult}{定义}{4.1}
设 $G$ 为 $S$-群。称一个带右 $G$ 作用的 $S$-\Term{prescheme}{预概形} $P$ 为 $G$ 下的（右）\Term{principal-homogeneous-bundle}{主齐性丛}，若存在 $S$ 的开\Term{covering}{覆盖} $(U_i)$，且对每个 $i$，存在\Term{faithfully-flat}{忠实平坦}且\Term{quasi-compact}{拟紧}的\Term{base-change}{基变换}\Term{morphism}{态射} $S_i'\to U_i$，使
\[
P'=P\times_SS'
\]
为 $G'=G\times_SS'$ 下\Term{trivial-bundle}{平凡}的带作用\Term{prescheme}{预概形}；这里 $S'$ 为各 $S_i'$ 的不交并所成的 $S$-\Term{prescheme}{预概形}。
\end{SGAresult}
（注意，\Term{base-change}{基变换}\Term{functor}{函子} $X\mapsto X'=X\times_SS'$ \Term{left-exact}{左正合}，因此将群变为群，将带\Term{group-action}{群作用}的对象变为带\Term{group-action}{群作用}的对象。）定义 4.1 在\Term{base-change}{基变换}下稳定。另有：

\begin{SGAresult}{命题}{4.2}
设 $G$ 为 $S$ 上\Term{flat}{平坦}且\Term{quasi-compact}{拟紧}的 $S$-群，$P$ 为带 $G$ 右作用的 $S$-\Term{prescheme}{预概形}。下列条件等价：
\SourcePage{243}
\begin{enumerate}[label=(\roman*)]
\item $P$ 为 $G$ 下的\Term{principal-homogeneous-bundle}{主齐性丛}；
\item $P$ 为 $G$ 下的\Term{formally-principal-homogeneous}{形式主齐性}对象，且结构\Term{morphism}{态射} $P\to S$ \Term{faithfully-flat}{忠实平坦}且\Term{quasi-compact}{拟紧}。
\end{enumerate}
\end{SGAresult}
若 $P$ 为 $G$ 下的\Term{principal-homogeneous-bundle}{主齐性丛}，沿用 4.1 的记号，$P'$ 在 $S'$ 上\Term{faithfully-flat}{忠实平坦}且\Term{quasi-compact}{拟紧}，因为 $G'$ 具有这些性质，且 $P'$ 与 $G'$ 在 $S'$ 上\Term{isomorphism}{同构}。因此 $P/S$ 也具有这些性质；“\Term{surjectivity}{满射}”和“\Term{quasi-compact}{拟紧}”见 VIII 3.1，而“\Term{flat}{平坦}”在第 VIII 讲的基本性质汇集中漏述了。反之，若满足（ii），取\Term{base-change}{基变换} $S'=P$，它确实在 $S$ 上\Term{faithfully-flat}{忠实平坦}且\Term{quasi-compact}{拟紧}。由于 $P$ 在 $S$ 上\Term{formally-principal-homogeneous}{形式主齐性}，而\Term{base-change}{基变换}\Term{functor}{函子}\Term{left-exact}{左正合}，$P'$ 在 $S'$ 上也\Term{formally-principal-homogeneous}{形式主齐性}。另一方面，$P'$ 在 $S'$ 上有一个\Term{section}{截面}，即对角\Term{section}{截面}，因此它是\Term{trivial-bundle}{平凡}\Term{principal-bundle}{主丛}。证明完毕。

\begin{SGAresult}{推论}{4.3}
若 $G$ 在 $S$ 上\Term{affine}{仿射}且\Term{flat}{平坦}，则任意 $G$ 下的\Term{principal-homogeneous-bundle}{主齐性丛} $P$ 都在 $S$ 上\Term{affine}{仿射}且\Term{flat}{平坦}。
\end{SGAresult}
事实上，作\Term{faithfully-flat}{忠实平坦}且\Term{quasi-compact}{拟紧}的\Term{base-extension}{基扩张}后，它具有这些性质，再应用（VIII 5.6）即可。

对于 $S$ 上\Term{flat}{平坦}\Term{affine}{仿射}的 $S$-群，定义 4.1 的作用在于（VIII 2.1），即：4.1 中的\Term{morphism}{态射} $S'\to S$ 是\Term{affine}{仿射}相对\Term{prescheme}{预概形}\Term{category}{范畴}的\Term{effective-descent-morphism}{有效下降态射}。凭借这一事实，下面概述的性质基本上都可用“\Term{category}{范畴}式”的方法验证。\footnote{参见脚注（3）所引文献（原书第 226 页）。}设 $E$ 为带 $S$-群 $G$ 左作用的 $S$-\Term{prescheme}{预概形}，$P$ 为 $G$ 下的右\Term{principal-homogeneous-bundle}{主齐性丛}。我们希望定义一个“局部”\Term{isomorphism}{同构}于 $E$ 的\Term{associated-bundle}{伴随丛} $E^{(P)}$。为此，令 $G$ 按照
\[
(x,y)\longmapsto(xg,g^{-1}y)
\]
右作用在 $P\times_SE$ 上；这个公式在集合情形描述该作用，并可用惯常的方法推广到\Term{category}{范畴}情形。在存在的前提下，置
\[
E^{(P)}=(P\times_SE)/G.
\]
于是 $P\times_SE$ 是 $T=E^{(P)}$ 上的\Term{prescheme}{预概形}，带有 $G_T=G\times_ST$ 的右作用。为使讨论顺畅，我们还希望 $P\times_SE$ 是 $T$ 上以 $G_T$ 为群的\Term{principal-homogeneous-bundle}{主齐性丛}。为证明 $E^{(P)}$ 的存在及这一性质，重新取定义 4.1 中的 $S'$，并考察原情形在 $S'$ 上的\Term{inverse-image}{逆像}。由于 $P'$ \Term{trivial-bundle}{平凡}，即\Term{isomorphism}{同构}于 $G_d'$，容易看出 $E'^{(P')}$ 存在且具有所需的\Term{exactness}{正合性}。事实上，$E'\times_{S'}P'$ 与 $E'\times_{S'}G'$ 为 $G'$-\Term{isomorphism}{同构}的，因而 $E'^{(P')}$ \Term{isomorphism}{同构}于 $E'$。此外，在\Term{trivial-bundle}{平凡}作用空间的情形，\Term{associated-bundle}{伴随丛}的构造与任意\Term{base-extension}{基扩张}交换。特别地，考虑\Term{base-extension}{基扩张}
\[
\begin{tikzcd}[column sep=large]
S''\arrow[r,shift left=.6ex]\arrow[r,shift right=.6ex]&S'
\end{tikzcd}
\qquad\text{及}\qquad
\begin{tikzcd}[column sep=large]
S'''\arrow[r,shift left=1.1ex]\arrow[r]\arrow[r,shift right=1.1ex]&S',
\end{tikzcd}
\]
其中 $S''$、$S'''$ 分别为 $S'$ 在 $S$ 上的二重及三重\Term{fiber-product}{纤维积}，
\SourcePage{244}
可见 $E'^{(P')}$ 带有相对于\Term{morphism}{态射} $S'\to S$ 的\Term{descent-datum}{下降资料}，且 $E^{(P)}$ 带有所需性质地存在，当且仅当该\Term{descent-datum}{下降资料}有效。此时 $E^{(P)}$ 当然就是\Term{descent}{下降}所得的对象。（这里使用 $S'\to S$ 是 $S$-\Term{prescheme}{预概形}\Term{category}{范畴}中的\Term{descent}{下降}\Term{morphism}{态射}，参见 VIII 5.2。）由此可知，若 $E$ 在 $S$ 上\Term{affine}{仿射}，\Term{associated-bundle}{伴随丛}便存在。

将此构造用于给定一个 $S$-\Term{group-homomorphism}{群同态} $G\to H$，并取 $E$ 为 $S$-\Term{prescheme}{预概形} $H$ 的情形，其中 $G$ 按给定\Term{homomorphism}{同态}在 $H$ 上作左作用。由于 $H$ 以右乘作用于自身，且此作用与 $G$ 的作用交换，而\Term{associated-bundle}{伴随丛}的构造在 $S$ 上存在时与\Term{base-extension}{基扩张}交换，容易看出 $H$ 将在 $P^{(H)}$ 上作右作用，使其成为 4.1 意义下 $H$ 下的\Term{principal-homogeneous-bundle}{主齐性丛}；更确切地说，它与 $P$ 可由同一个\Term{morphism}{态射} $S'\to S$ \Term{trivial-bundle}{平凡}化。因此，对 $G$ 下任意\Term{principal-homogeneous-bundle}{主齐性丛} $P$ 以及任意 $S$-\Term{group-homomorphism}{群同态} $G\to H$，若 $H$ 在 $S$ 上\Term{affine}{仿射}，则相应地得到一个以 $H$ 为群的\Term{principal-homogeneous-bundle}{主齐性丛}。该构造对 $(G\to H)$ 具有\Term{functorial}{函子性}，并与任意\Term{base-change}{基变换} $T\to S$ 相容。

\begin{SGAresult}{定义}{4.4}
设 $G$ 为 $S$-\Term{prescheme}{预概形}。以 $H^0(S,G)$ 表示 $G$ 在 $S$ 上的\Term{section}{截面}集；当 $G$ 为 $S$-群时，将其视为群。此时，以 $H^1(S,G)$ 表示 $S$ 上群为 $S$ 的\Term{principal-homogeneous-bundle}{主齐性丛}的\Term{isomorphism}{同构}类集合，并在其中以\Term{trivial-bundle}{平凡}丛所对应的元素为“标记点”。\footnote{只有在具备\Term{descent}{下降}有效性判据时，这一记号才与一般\Term{cohomology}{上同调}记号（SGA 4 V）一致；通常只有对\Term{affine}{仿射} $G$，或仅拟\Term{affine}{仿射}的 $G$（参见 VIII 7.9），才有这样的保证。}
\end{SGAresult}
因此，$H^0(S,G)$ 是关于 $S$-\Term{prescheme}{预概形} $G$ 的取值于集合\Term{category}{范畴}的\Term{functor}{函子}。它\Term{left-exact}{左正合}，因而与\Term{finite}{有限}乘积交换；这就意味着它将群变为群，将交换群变为交换群。同样，借助\Term{associated-bundle}{伴随丛}的构造，$H^1(S,G)$ 是关于\Term{affine}{仿射} $S$-群 $G$ 的集合值\Term{functor}{函子}；很容易验证它与\Term{finite}{有限}乘积交换。特别地，它将\Term{affine}{仿射} $S$-群\Term{category}{范畴}中的群，即\Term{affine}{仿射}交换 $S$-群，变为后一\Term{category}{范畴}中的群，而且是交换群，因为前一\Term{category}{范畴}中的群是交换的。因此，若 $G$ 为\Term{affine}{仿射}交换 $S$-群，则 $H^1(S,G)$ 为交换群；\Term{affine}{仿射}交换 $S$-\Term{group-homomorphism}{群同态} $G\to H$ 则给出\Term{group-homomorphism}{群同态} $H^1(S,G)\to H^1(S,H)$。

为简便，以下只考虑\Term{affine}{仿射}交换 $S$-群。设
\[
0\longrightarrow G'\xrightarrow{u}G\xrightarrow{v}G''\longrightarrow0
\]
为这样的群之间的一列\Term{morphism}{态射}。
\SourcePage{245}
若 $vu=0$，因而可将 $G$ 视为带 $G'_{G''}$ 右作用的 $G''$ 上\Term{prescheme}{预概形}，且 $G$ 是 $G''$ 上以 $G'_{G''}=G'\times_SG''$ 为群的\Term{principal-homogeneous-bundle}{主齐性丛}，则称此列\Term{exact-sequence}{正合}。这特别蕴含 $u:G'\to G$ 为 $v$ 的核，因而还蕴含下列序列\Term{exact}{正合}：
\[
0\longrightarrow H^0(X,G')\longrightarrow H^0(X,G)\longrightarrow H^0(X,G'').
\]
此外，还可以定义映射
\[
\partial:H^0(X,G'')\longrightarrow H^1(X,G'):
\]
对 $G''$ 在 $S$ 上的任意\Term{section}{截面}，即 $S$-\Term{morphism}{态射} $f:S\to G''$，赋予 $G''$ 上的\Term{principal-homogeneous-bundle}{主齐性丛} $G$ 的\Term{inverse-image}{逆像} $P_f$，它是 $S$ 上以 $G'\simeq f^*(G'_{G''})$ 为群的\Term{principal-homogeneous-bundle}{主齐性丛}。从 $S$-\Term{prescheme}{预概形}的观点看，它正是 $S$ 经\Term{immersion}{浸入} $f$ 所得像子\Term{prescheme}{预概形}在 $v:G\to G''$ 下的\Term{inverse-image}{逆像}；$P_f$ 上的 $G'$ 作用由 $G'$ 在 $G$ 上的右作用诱导。下面命题的验证也留给读者；除了书写上的繁琐，并无其他困难。

\begin{SGAresult}{命题}{4.5}
映射 $\partial:H^0(X,G'')\to H^1(X,G')$ 为\Term{group-homomorphism}{群同态}，且下列\Term{homomorphism}{同态}序列\Term{exact}{正合}：
\[
\begin{aligned}
0\longrightarrow H^0(X,G')&\xrightarrow{u^0}H^0(X,G)
\xrightarrow{v^0}H^0(X,G')\xrightarrow{\partial}H^1(X,G'')\\
&\xrightarrow{u^1}H^1(X,G)\xrightarrow{v^1}H^1(X,G'').
\end{aligned}
\]
这里，除 $\partial$ 外的\Term{homomorphism}{同态}均来自 $H^0$、$H^1$ 的\Term{functor}{函子}结构。
\end{SGAresult}

\textbf{注记 4.6。} 这里研究\Term{principal-homogeneous-bundle}{主齐性丛}的观点显然受 Serre [7] 的启发，建议读者参阅该文。若希望得到同样适用于 $S$ 上拟\Term{projective}{射影}结构群的形式体系，尤其要包含\Term{projective}{射影}阿贝尔\Term{scheme}{概形}，则可修改 4.1，要求 $S'$ 为\Term{prescheme}{预概形} $S_i'$ 的和，其中各 $S_i'$ 在\Term{covering}{覆盖} $S$ 的开集 $S_i$ 上\Term{finite}{有限}且\Term{locally-free}{局部自由}。此时，将上面所有“\Term{affine}{仿射}”假设换为“拟\Term{projective}{射影}”，并相应解释此前所给的 $S$-群\Term{exact-sequence}{正合列}定义，前述论述，特别是 4.5，都仍然成立。事实上，只须将对（VIII 2.1）的引用改为（VIII 7.7）：所用\Term{morphism}{态射} $S'\to S$ 是相对拟\Term{projective}{射影}\Term{prescheme}{预概形}的\Term{fibered-category}{纤维范畴}中的\Term{effective-descent-morphism}{有效下降态射}。但须注意，这第二种\Term{principal-homogeneous-bundle}{主齐性丛}的概念比 4.1 的第一种概念更严格。

\SourcePage{246}
\textbf{4.7。} 若要求 $S$ 可由开集 $S_i$ \Term{covering}{覆盖}，使对每个 $i$，$P|_{S_i}$ 为 $G|_{S_i}$ 下的\Term{trivial-bundle}{平凡}作用丛，则得到更严格的\Term{principal-homogeneous-bundle}{主齐性丛}概念；此时称 $P$ 为\Term{locally-trivial-principal-bundle}{局部平凡主齐性丛}。给定 $G$ 后，这些丛的\Term{isomorphism}{同构}类构成 $H^1(X,G)$ 的一个子集，它与 $H^1(X,\mathscr O_X(G))$ 一一对应，其中 $\mathscr O_X(G)$ 为 $G$ 在 $S$ 上的\Term{section}{截面}所成的通常意义下的层，参见 [2]。若要这些 $H^1$ 仍给出 4.5 那样的\Term{cohomology}{上同调}\Term{exact-sequence}{正合列}，显然须假定 $0\to G'\to G\to G''\to0$ 在这一新情境下以合理的意义\Term{exact}{正合}，即 $G$ 是 $G''$ 上以 $G'_{G''}$ 为群的\Term{locally-trivial-bundle}{局部平凡丛}。等价地，$u:G'\to G$ 为 $v:G\to G''$ 的核，且 $G$ 在 $G''$ 上局部有\Term{section}{截面}。

\textbf{4.8。} 自然希望引入高阶\Term{cohomology}{上同调}群 $H^i(X,G)$，以继续延伸\Term{exact-sequence}{正合列} 4.5。这可以在“Weil \Term{cohomology}{上同调}”的框架下做到：考虑 $S$ 上\Term{quasi-compact}{拟紧}\Term{prescheme}{预概形}的\Term{category}{范畴} $\mathcal B$，并指定其中\Term{faithfully-flat}{忠实平坦}且\Term{quasi-compact}{拟紧}\Term{morphism}{态射}的集合 $\mathcal M$，称之为\Term{localization}{局部化}\Term{morphism}{态射}。$S$ 上，或更确切地说 $(\mathcal B,\mathcal M)$ 上的阿贝尔“\Term{weil-sheaf}{Weil 层}”，是一个从 $\mathcal B$ 到\Term{abelian-group}{阿贝尔群}\Term{category}{范畴}的\Term{contravariant-functor}{反变函子} $\mathcal F$，它将和变为积，并将 $f\in\mathcal M$ 的序列
\[
\begin{tikzcd}[column sep=large]
T''=T'\times_TT'\arrow[r,shift left=.6ex,"\operatorname{pr}_1"]\arrow[r,shift right=.6ex,"\operatorname{pr}_2"']&T'\arrow[r,"f"]&T
\end{tikzcd}
\]
变为\Term{exact}{正合}的集合图表
\[
\begin{tikzcd}[column sep=large]
\mathcal F(T)\arrow[r]&\mathcal F(T')\arrow[r,shift left=.6ex]\arrow[r,shift right=.6ex]&\mathcal F(T'').
\end{tikzcd}
\]
\Term{weil-sheaf}{Weil 层}构成一个\Term{inductive-limit}{归纳极限}\Term{exact}{正合}且具有生成元的\Term{abelian-category}{阿贝尔范畴}，因而具有足够多的\Term{injective-object}{内射对象} [1]。\Term{functor}{函子} $\Gamma(\mathcal F)=\mathcal F(S)$ 的\Term{right-derived-functor}{右导出函子}记为 $H^i(S,\mathcal F)$。另一方面，每个交换 $S$-群显然都定义一个 \Term{weil-sheaf}{Weil 层}（VIII 5.2），其 $H^0$、$H^1$ 正是 $H^0(S,G)$、$H^1(S,G)$，从而可合理地定义其余的 $H^i(S,G)$。还可证明，$S$-群的\Term{exact-sequence}{正合列}定义 \Term{weil-sheaf}{Weil 层}的\Term{exact-sequence}{正合列}，因此可以重新得到并延伸\Term{exact-sequence}{正合列} 4.5。\footnote{对这一观点的系统研究，见 SGA 4 I—IX。}

\textbf{4.9。} 如 [2] 那样发展 4.5 的非交换变体，是很合适的。关于在适当框架下对本节所概述的各种\Term{cohomology}{上同调}概念的系统论述，请参阅 J. Giraud 正在准备的工作。\footnote{参见 J. Giraud，\textit{Cohomologie non abélienne}，Springer-Verlag，1971。}

\SGAsection{5}{主丛的几个特殊情形}
现在假设 $S$ \Term{connected}{连通}，并给定一个\Term{geometric-point}{几何点} $a$，从而有\Term{fundamental-group}{基本群} $\pi_1(S,a)$，它分类 $S$ 的\Term{etale-covering}{平展覆盖}：
\SourcePage{247}
$S$ 的\Term{etale-covering}{平展覆盖}\Term{category-equivalence}{范畴等价}于带 $\pi_1$ 连续作用的\Term{finite}{有限}集合\Term{category}{范畴}。因此，$S$ 上\Term{finite}{有限}且\Term{etale}{平展}的\Term{group-scheme}{群概形} $G$ 本质上由一个通常的\Term{finite-group}{有限群} $\mathscr G$ 确定，$\pi_1$ 以群\Term{automorphism}{自同构}连续作用于它。$S$ 的带 $G$ 右作用的\Term{etale-covering}{平展覆盖} $P$，本质上由一个带 $\pi_1$ 连续左作用的\Term{finite}{有限}集合 $\mathscr P$ 确定，其上还有与 $\pi_1$ 作用相容的 $\mathscr G$ 右作用：
\[
s(p\cdot g)=(sp)\cdot(sg)\qquad
(s\in\pi_1,\ p\in\mathscr P,\ g\in\mathscr G).
\]
可验证，$P$ 为 4.1 意义下的\Term{principal-homogeneous-bundle}{主齐性丛}，当且仅当 $\mathscr P$ 为 $\mathscr G$ 下的\Term{principal-homogeneous}{主齐性}集合，例如可使用判据 4.2。换言之，$S$ 上以 $G$ 为群的\Term{principal-homogeneous-bundle}{主齐性丛}的\Term{category}{范畴}，等价于带 $\pi_1$ 连续作用的\Term{finite}{有限}集合\Term{category}{范畴}中以 $\mathscr G$ 为群的\Term{principal-homogeneous-bundle}{主齐性丛}的\Term{category}{范畴}。特别地，得到一个对 $G$ 具有\Term{functorial}{函子性}的典范\Term{bijection}{双射}
\[
H^1(S,G)\xrightarrow{\sim}H^1(\pi_1,\mathscr G).\tag{*}
\]
右端表示带 $\pi_1$ 作用的\Term{finite}{有限}集合\Term{category}{范畴}中，$\mathscr G$ 下\Term{principal-homogeneous-bundle}{主齐性丛}的\Term{isomorphism}{同构}类集合；这里其实无须再特别说明作用连续。按熟知的方法，它可明确写成 1-\Term{cocycle}{上循环} $\varphi:\pi_1\to\mathscr G$ 的集合 $Z^1(\pi_1,\mathscr G)$ 对群 $\mathscr G$ 的自然作用所取的商，其中
\[
\varphi(1)=1,\qquad \varphi(st)=\varphi(s)\bigl(s\cdot\varphi(t)\bigr).
\]
一个重要情形是 $\pi_1$ 在 $\mathscr G$ 上的作用\Term{trivial-bundle}{平凡}，即 $G$ 是 $S$ 的\Term{completely-split}{完全分裂}\Term{covering}{覆盖}，\Term{isomorphism}{同构}于 $\mathscr G$ 份 $S$ 的和。此时也以 $H^1(S,\mathscr G)$ 代替 $H^1(S,G)$；由（*），此集合与
\[
H^1(\pi_1,\mathscr G)
=\operatorname{Hom}(\pi_1,\mathscr G)/\mathscr G\text{ 的内自同构}
\]
一一对应。还应注意，此时 $S$ 上以 $G$ 为群的\Term{principal-homogeneous-bundle}{主齐性丛}，无非就是 $S$ 的以 $\mathscr G$ 为群的\Term{principal-covering}{主覆盖}（V 2.7）；上述\Term{bijection}{双射}也可由下述对应得到：在 $a$ 上方带点的、以 $\mathscr G$ 为群的 $S$ 的\Term{principal-covering}{主覆盖}，与 $\pi_1(S,a)$ 到 $\mathscr G$ 的连续\Term{homomorphism}{同态}相对应（第 V 讲第 5 节末）。

将\Term{etale-covering}{平展覆盖}理论与\Term{principal-bundle}{主丛}理论相联系的重要性——这一联系已隐含于 A. Weil 的《Généralisation des Fonctions Abéliennes》，并由 S. Lang 在其几何类域论中明确提出，参见 Serre [5]——在于：$S$ 上任意\Term{finite}{有限}\Term{etale}{平展} $S$-群，都可嵌入一个 $S$-群 $H$，后者在 $S$ 上\Term{affine}{仿射}且\Term{smooth}{光滑}，\Term{fiber}{纤维}\Term{connected}{连通}，并在 $G$ 交换时也交换。因此，通过\Term{exact-sequence}{正合列} 4.5，必要时使用其非交换变体，可以借助以 $H$ 为群的\Term{principal-bundle}{主丛}的“连续”分类，研究以 $G$ 为群的\Term{principal-covering}{主覆盖}的“\Term{discrete}{离散}”分类；反之亦然。关于一般嵌入 $G\to H$ 的构造思想，参见 [5，VI 2.8]，不过这一构造似乎很少在实践中使用。
\SourcePage{248}
下一节只展开两个重要且经典的特殊情形。为此需要一个辅助结果：

\begin{SGAresult}{命题}{5.1}
设 $S$ 为\Term{prescheme}{预概形}，$G$ 为与 $\operatorname{GL}(n)_S$（例如 $\mathbf G_{m,S}$）或 $\mathbf G_{a,S}$ \Term{isomorphism}{同构}的 $S$-群。则 $G$ 下的任意\Term{principal-homogeneous-bundle}{主齐性丛}均局部\Term{trivial-bundle}{平凡}。
\end{SGAresult}
具体地，$\operatorname{GL}(n)_S$（$n\geq0$ 为整数）表示代表下列关于 $S$-\Term{prescheme}{预概形} $T$ 的\Term{contravariant-functor}{反变函子}的 $S$-群：
\[
T\longmapsto\operatorname{GL}\bigl(n,\Gamma(T,O_T)\bigr).
\]
特别地，$\mathbf G_{m,S}$（“$S$ 上的\Term{multiplicative-group}{乘法群}”）代表\Term{contravariant-functor}{反变函子} $T\mapsto\Gamma(T,O_T^*)$，因此作为 $S$ 上的\Term{prescheme}{预概形}，它\Term{isomorphism}{同构}于 $\operatorname{Spec}O_S[t,t^{-1}]$，其中 $t$ 为不定元。同样，$\mathbf G_{a,S}$ 代表\Term{contravariant-functor}{反变函子} $T\mapsto\Gamma(T,O_T)$，因而作为 $S$-\Term{prescheme}{预概形}\Term{isomorphism}{同构}于 $\operatorname{Spec}(O_S[t])$。注意，通过\Term{devissage}{分解论证}，5.1 重新给出 Rosenlicht 的局部\Term{trivial-bundle}{平凡}性结果：若 $G$ 具有一个“合成列”，其相邻因子均为这里讨论的群，则也有局部\Term{trivial-bundle}{平凡}性。关于\Term{principal-homogeneous-bundle}{主齐性丛}局部\Term{trivial-bundle}{平凡}性问题的更精细研究，参见 [7] 和 [3]。

第一个断言的证明只须注意 $G(T)=\operatorname{Aut}(O_T^n)$，并且 4.1 中出现的\Term{faithfully-flat}{忠实平坦}\Term{quasi-compact}{拟紧}\Term{morphism}{态射} $S'\to S$，是\Term{local-isomorphism}{局部同构}于 $O_T^n$ 的模、即秩 $n$ 的\Term{locally-free}{局部自由}模所成\Term{fibered-category}{纤维范畴}的\Term{effective-descent-morphism}{有效下降态射}（VIII 1.12）。第二个断言类似：此时 $G(T)=\operatorname{Aut}(E_T)$，其中 $E_T$ 为 $O_T$ 被 $O_T$ 的\Term{trivial-bundle}{平凡}\Term{extension}{扩张}，而\Term{automorphism}{自同构}当然须保持\Term{extension}{扩张}结构。4.1 中的\Term{morphism}{态射} $S'\to S$，是 $O_T$ 被 $O_T$ 的\Term{extension}{扩张}所成\Term{fibered-category}{纤维范畴}的\Term{effective-descent-morphism}{有效下降态射}；这很容易由 VIII 1.1 得到。而这样的\Term{extension}{扩张}自动局部\Term{trivial-bundle}{平凡}。

\textbf{注记 5.2。} 相同类型的证明适用于\Term{symplectic-group}{辛群} $\operatorname{Sp}(2n)_S$，因为在一个\Term{local-isomorphism}{局部同构}于 $O_S^{2n}$ 的模上的交错形式，若“非退化”，即给出该模到其对偶的\Term{isomorphism}{同构}，则\Term{local-isomorphism}{局部同构}于标准形式。反之，对正交群结论不成立；即使 $S$ 为域的谱也不成立，因为域上的二次形式可能不与标准形式\Term{isomorphism}{同构}。事实上，[3] 中基本上证明了：$\operatorname{GL}$、$\operatorname{Sp}$、$\mathbf G_a$ 及那些\Term{separable}{可分}解为这类群的群，差不多是仅有的使这里这类局部\Term{trivial-bundle}{平凡}性结果成立的群。

\begin{SGAresult}{推论}{5.3}
有典范\Term{bijection}{双射}
\[
H^1(S,\operatorname{GL}(n)_S)\xleftarrow{\sim}H^1(S,\operatorname{GL}(n,O_S)),
\]
\SourcePage{249}
特别地，
\[
H^1(S,\mathbf G_{m,S})\xleftarrow{\sim}H^1(S,O_S^*),
\]
以及
\[
H^1(S,\mathbf G_{a,S})\xleftarrow{\sim}H^1(S,O_S),
\]
其中各右端均为\Term{topological-space}{拓扑空间} $S$ 以通常的层为系数的\Term{cohomology}{上同调}。
\end{SGAresult}
特别地，$H^1(S,\operatorname{GL}(n)_S)$ 可认同于 $S$ 上秩 $n$ 的\Term{locally-free}{局部自由}模的\Term{isomorphism}{同构}类集合，而 $H^1(S,\mathbf G_{a,S})$ 可认同于模 $O_S$ 被自身的\Term{extension}{扩张}的等价类集合。

\SGAsection{6}{应用于主覆盖：Kummer 理论与 Artin–Schreier 理论}
\begin{SGAresult}{命题}{6.1}
设 $S$ 为\Term{prescheme}{预概形}，$n>0$ 为整数，$u_n:\mathbf G_{m,S}\to\mathbf G_{m,S}$ 为取 $n$ 次幂的\Term{homomorphism}{同态}，$\mu_{n,S}$ 为其核。则 $\mu_{n,S}$ 在 $S$ 上\Term{finite}{有限}且\Term{locally-free}{局部自由}，秩为 $n$；并且它在 $S$ 上\Term{etale}{平展}，当且仅当对每个 $s\in S$，$s$ 的特征都与 $n$ 互素。\Term{homomorphism}{同态}序列
\[
0\longrightarrow\mu_{n,S}\longrightarrow\mathbf G_{m,S}
\xrightarrow{u_n}\mathbf G_{m,S}\longrightarrow0
\]
在第 4 节的意义下\Term{exact}{正合}，称为 $S$ 上对应于整数 $n$ 的\Term{kummer-exact-sequence}{Kummer 正合列}。
\end{SGAresult}
有
\[
\mathbf G_m=\operatorname{Spec}O_S[t,t^{-1}],
\]
而 $u_n$ 对应于\Term{affine}{仿射} $O_S$-代数上的\Term{homomorphism}{同态} $u_n$，其定义为
\[
u_n(t)=t^n.
\]
另一方面，$\mathbf G_{m,S}$ 的单位\Term{section}{截面}对应于下列 $O_S$-代数的增广\Term{homomorphism}{同态}：
\[
\varepsilon(t)=1,
\]
其核为主理想 $(t-1)$。该理想在 $u_n$ 下的像为主理想 $(1-t^n)$，因而得到
\[
\mu_{n,S}=\operatorname{Spec}O_S[t]/(1-t^n).
\]
这特别表明，$\mu_{n,S}$ 在 $S$ 上\Term{finite}{有限}，且由 $S$ 上秩 $n$ 的自由代数定义，其基为 $t^i$（$0\leq i\leq n-1$）。它在 $s\in S$ 处\Term{etale}{平展}，当且仅当 $k[t]/(1-t^n)$ 在 $k=k(s)$ 上\Term{separable}{可分}；该代数由向 $k$ 形式添入所有 $n$ 次单位根得到。等价地，$1-t^n$ 在 $k$ 的一个\Term{algebraic-closure}{代数闭包}中的根两两不同，而这等价于 $n$ 与特征互素。最后，为证明 6.1 的\Term{homomorphism}{同态}序列\Term{exact}{正合}，由判据 4.2，只须证明 $u_n$ \Term{faithfully-flat}{忠实平坦}。
\SourcePage{250}
显然可假定 $S$ \Term{affine}{仿射}，其环为 $A$。于是 $\mathbf G_{m,S}$ \Term{affine}{仿射}，环为 $B=A[t,t^{-1}]$；只须验证 $u_n$ 使 $B$ 成为 $B$ 上秩 $n$ 的\Term{free-module}{自由模}。等价地，须证 $u_n$ 为\Term{injectivity}{单射}，且 $A[t,t^{-1}]$ 是 $A[t^n,t^{-n}]$ 上秩 $n$ 的\Term{free-module}{自由模}。事实上，容易验证 $t^i$（$0\leq i\leq n-1$）构成这样的一组基。证明完毕。

\begin{SGAresult}{定义}{6.2}
称 $\mu_{n,S}$ 为 $S$ 上秩 $n$ 的\Term{kummer-group}{Kummer 群}；称 $S$ 上以此群为结构群的任意\Term{principal-homogeneous-bundle}{主齐性丛}为 $S$ 上秩 $n$ 的\Term{kummer-principal-covering}{Kummer 主覆盖}。
\end{SGAresult}
这些\Term{covering}{覆盖}的集合构成一个群，记为 $H^1(S,\mu_{n,S})$，或简记为 $H^1(S,\mu_n)$。注意，$S$ 上秩 $n$ 的 \Term{kummer-group}{Kummer 群}的构造与\Term{base-extension}{基扩张}相容，因此 $\mu_{n,S}$ 来自 $\operatorname{Spec}(\mathbf Z)$ 上的绝对 \Term{kummer-group}{Kummer 群} $\mu_n$ 的\Term{base-extension}{基扩张}。

以 $(\mathbf Z/n\mathbf Z)_S$ 表示由通常\Term{finite-group}{有限群} $\mathbf Z/n\mathbf Z$ 定义的 $S$-群。若 $G$ 为任意 $S$-群，则由 $(\mathbf Z/n\mathbf Z)_S$ 到 $G$ 的 $S$-\Term{group-homomorphism}{群同态} $u$，与 $G$ 在 $S$ 上那些 $n$ 次幂为单位\Term{section}{截面}的\Term{section}{截面}一一对应，且这一对应与\Term{base-change}{基变换}相容：赋予 $u$ 的\Term{section}{截面}，正是 $(\mathbf Z/n\mathbf Z)_S$ 在 $S$ 上由 $\mathbf Z/n\mathbf Z$ 的生成元 $1\bmod n\mathbf Z$ 定义的\Term{section}{截面}在 $u$ 下的像。据此有：

\begin{SGAresult}{推论}{6.3}
若 $\mu_{n,S}$ 在 $S$ 上\Term{etale}{平展}，则由上述方法，$S$-群\Term{isomorphism}{同构}
\[
(\mathbf Z/n\mathbf Z)_S\xrightarrow{\sim}\mu_{n,S}
\]
与 $O_S$ 的那些在 $S$ 的每个\Term{connected-component}{连通分支}上恰为 $n$ 阶的\Term{section}{截面}一一对应；这样的\Term{section}{截面}称为“$S$ 上的\Term{primitive-root-of-unity}{本原 $n$ 次单位根}”。因此，$\mu_{n,S}$ 作为 $S$-群与 $(\mathbf Z/n\mathbf Z)_S$ \Term{isomorphism}{同构}，当且仅当它在 $S$ 上\Term{etale}{平展}，即 $S$ 的各\Term{residue-characteristic}{剩余特征}与 $n$ 互素，且 $S$ 上存在本原 $n$ 次单位根。
\end{SGAresult}
这解释了在经典 Kummer 理论中，为何要求扮演 $S$ 角色的基域的特征与 $n$ 互素，并包含 $n$ 次单位根，以及为何要选定一个本原 $n$ 次单位根。有了\Term{scheme}{概形}语言，就无须再受这些假设束缚；应当直接使用 $\mu_n$，而不是 $\mathbf Z/n\mathbf Z$。例如，将 6.1、4.5 和 5.3 结合起来，便得到 \Term{kummer-principal-covering}{Kummer 主覆盖}理论与 \Term{picard-group}{Picard 群}理论之间的以下一般关系：

\begin{SGAresult}{命题}{6.4}
设 $S$ 为\Term{prescheme}{预概形}，则有典范\Term{exact-sequence}{正合列}
\[
\begin{aligned}
0\longrightarrow H^0(S,\mu_n)&\longrightarrow H^0(S,O_S^*)
\longrightarrow H^0(S,O_S^*)\\
&\longrightarrow H^1(S,\mu_n)\longrightarrow H^1(S,O_S^*)
\longrightarrow H^1(S,O_S^*).
\end{aligned}
\]
\SourcePage{251}
因此，置 $H^1(S,O_S^*)=\operatorname{Pic}(S)$，并对任意\Term{abelian-group}{阿贝尔群} $A$，分别以 $ {}_nA$、$A_n$ 表示 $A$ 上乘 $n$ \Term{homomorphism}{同态}的核与\Term{cokernel}{余核}，可得\Term{exact-sequence}{正合列}
\[
0\longrightarrow H^0(S,O_S)^*_n
\longrightarrow H^1(S,\mu_n)\longrightarrow{}_n\operatorname{Pic}(S)
\longrightarrow0.
\]
\end{SGAresult}
下面具体说明两个重要情形，其中此\Term{exact-sequence}{正合列}的某一个端项为零。

\begin{SGAresult}{推论}{6.5}
假设 $ {}_n\operatorname{Pic}(S)=0$，例如 $S$ 为某个\Term{local-ring}{局部环}或阶乘环的谱；并设 $A=H^0(S,O_S)$。则有典范\Term{isomorphism}{同构}
\[
H^1(S,\mu_n)\xrightarrow{\sim}A^*/A^{*n}.
\]
\end{SGAresult}
当 $S$ 为域的谱时，这基本上就是经典 Kummer 理论的表述。

\begin{SGAresult}{推论}{6.6}
假设 $H^0(S,O_S)$ 的每个元素均为 $n$ 次幂，例如 $H^0(S,O_S)$ 为若干\Term{algebraically-closed-field}{代数闭域}的直积，或者 $S$ 在\Term{algebraically-closed-field}{代数闭域} $k$ 上\Term{reduced}{约化}且\Term{proper}{固有}。则有典范\Term{isomorphism}{同构}
\[
H^1(S,\mu_n)\xrightarrow{\sim}{}_n\operatorname{Pic}(S).
\]
\end{SGAresult}
特别地，当 $S$ 在\Term{algebraically-closed-field}{代数闭域} $k$ 上\Term{proper}{固有}且\Term{connected}{连通}时，这就将 $S$ 的\Term{fundamental-group}{基本群}与 $S/k$ 的\Term{picard-scheme}{Picard 概形} $P$ 的\Term{finite}{有限}阶点联系起来；例如，对与特征互素的 $n$，有\Term{isomorphism}{同构}
\[
\operatorname{Hom}(\pi_1(S),\mathbf Z/n\mathbf Z)\simeq{}_nP(k),
\]
这在代数几何中经常使用。作为应用，若 $P$ 的\Term{connected-component}{连通分支} $P^0$ 为 $g$ 维完备\Term{group-scheme}{群概形}，利用第 2 节所回忆的结果以及 Néron–Severi 扭群的\Term{finite}{有限}性，可以看出：对每个与特征互素的素数 $\ell$，\Term{fundamental-group}{基本群} $\pi_1(S)$ 的阿贝尔化的 $\ell$-准素分量，是 $\ell$ 进整数环 $\mathbf Z_\ell$ 上秩 $2g$ 的\Term{finite-type}{有限型}模；而且，除至多\Term{finite}{有限}多个 $\ell$ 外，该模自由。如 Serre 所注意到的，这允许我们在一定条件下，利用\Term{semicontinuity}{半连续性}定理（X 2.3）证明：若 $X$ 为\Term{connected}{连通}\Term{scheme}{概形} $S$ 上的\Term{flat}{平坦}\Term{projective}{射影}\Term{scheme}{概形}，则 $X$ 的各\Term{fiber}{纤维}的 \Term{picard-scheme}{Picard 概形}都具有同一维数。只要 $X/S$ 的 \Term{picard-scheme}{Picard 概形}存在，且 $X/S$ 各\Term{fiber}{纤维}的\Term{connected}{连通} \Term{picard-scheme}{Picard 概形}都是\Term{proper}{固有}\Term{group-scheme}{群概形}，Serre 的论证就适用。例如，当 $X/S$ 的\Term{geometric-fiber}{几何纤维}\Term{normal}{正规}时就如此；这里仍假定 $X/S$ \Term{flat}{平坦}且\Term{projective}{射影}。特别地，$X/S$ \Term{smooth}{光滑}且\Term{projective}{射影}时适用。

现在设 $p$ 为素数，且 $S$ 为特征 $p$ 的\Term{prescheme}{预概形}，即 $p\cdot O_S=0$。于是 $O_S$ 中取 $p$ 次幂的映射是加性的；
\SourcePage{252}
将 $S$ 换为在 $S$ 上变化的 $T$，得到的相应\Term{morphism}{态射}
\[
F:\mathbf G_{a,S}\longrightarrow\mathbf G_{a,S}
\]
因而是 $S$-\Term{group-homomorphism}{群同态}，称为\Term{frobenius-homomorphism}{Frobenius 同态}。（注意：对任何由素域 $\mathbf Z/p\mathbf Z$ 上的\Term{prescheme}{预概形} $G_0$ \Term{base-extension}{基扩张}得到的 $S$-\Term{prescheme}{预概形} $G$，都可定义这样的\Term{morphism}{态射}；若 $G_0$ 为\Term{group-prescheme}{群预概形}，该\Term{morphism}{态射}就是\Term{group-homomorphism}{群同态}。）置
\[
\wp=\operatorname{id}-F:\mathbf G_{a,S}\longrightarrow\mathbf G_{a,S}.
\]
另一方面，考虑由通常\Term{finite-group}{有限群} $\mathbf Z/p\mathbf Z$ 定义的 $S$-群 $(\mathbf Z/p\mathbf Z)_S$。前面已说过：对任意 $S$-群 $G$，从 $(\mathbf Z/p\mathbf Z)_S$ 到 $G$ 的 $S$-\Term{group-homomorphism}{群同态}，与 $G$ 在 $S$ 上那些 $p$ 次幂为单位\Term{section}{截面}的\Term{section}{截面}一一对应。当 $G=\mathbf G_{a,S}$ 时，它们便对应于 $G$ 在 $S$ 上的任意\Term{section}{截面}。特别地，取 $\mathbf G_{a,S}$ 在 $S$ 上与环层 $O_S$ 的单位\Term{section}{截面}对应的\Term{section}{截面}，得到 $S$-\Term{group-homomorphism}{群同态}
\[
i:(\mathbf Z/p\mathbf Z)_S\longrightarrow\mathbf G_{a,S}.
\]

\begin{SGAresult}{命题}{6.7}
$S$-\Term{group-homomorphism}{群同态}序列
\[
0\longrightarrow(\mathbf Z/p\mathbf Z)_S
\longrightarrow\mathbf G_{a,S}\longrightarrow\mathbf G_{a,S}\longrightarrow0
\]
在第 4 节的意义下\Term{exact}{正合}，称为 $S$ 上的\Term{artin-schreier-exact-sequence}{Artin–Schreier 正合列}。
\end{SGAresult}
只须在素域 $k=\mathbf Z/p\mathbf Z$ 上证明。注意到由 $\wp^*(t)=t-t^p$ 定义的\Term{homomorphism}{同态} $\wp^*:k[t]\to k[t]$ 使 $k[t]$ 成为 $k[t]$ 上秩 $p$ 的\Term{free-module}{自由模}即可；更确切地，若 $s=t-t^p$，则 $k[t]$ 为 $k[s]$ 上的\Term{free-module}{自由模}，以 $t^i$（$0\leq i\leq p-1$）为一组基。

由此结合 4.5 和 5.3，得到：

\begin{SGAresult}{命题}{6.8}
有典范\Term{exact-sequence}{正合列}
\[
\begin{aligned}
0\longrightarrow H^0(S,\mathbf Z/p\mathbf Z)&\longrightarrow H^0(S,O_S)
\longrightarrow H^0(S,O_S)\\
&\longrightarrow H^1(S,\mathbf Z/p\mathbf Z)
\longrightarrow H^1(S,O_S)\longrightarrow H^1(S,O_S),
\end{aligned}
\]
从而得到\Term{exact-sequence}{正合列}
\[
0\longrightarrow H^0(S,O_S)/\wp H^0(S,O_S)
\longrightarrow H^1(S,\mathbf Z/p\mathbf Z)
\longrightarrow H^1(S,O_S)^F\longrightarrow0.
\]
最后一项中的上标 $F$ 表示自\Term{homomorphism}{同态} $F$ 的不变量子群，即 $\wp=\operatorname{id}-F$ 的核。
\end{SGAresult}
再来说明两个极端情形。
\SourcePage{253}
\begin{SGAresult}{推论}{6.9}
假设 $H^1(S,O_S)^F=0$，例如 $S$ 为\Term{affine}{仿射}\Term{scheme}{概形}。置 $A=H^0(S,O_S)$，则有典范\Term{isomorphism}{同构}
\[
H^1(S,\mathbf Z/p\mathbf Z)\xrightarrow{\sim}A/\wp A.
\]
\end{SGAresult}
这就是经典形式的 Artin–Schreier 理论，至少在 $A$ 为域的谱时如此。

\begin{SGAresult}{推论}{6.10}
假设 $\wp H^0(S,O_S)=H^0(S,O_S)$，例如 $H^0(S,O_S)$ 为若干\Term{algebraically-closed-field}{代数闭域}的直积，或者 $S$ 在\Term{algebraically-closed-field}{代数闭域}上\Term{proper}{固有}。则有典范\Term{isomorphism}{同构}
\[
H^1(S,\mathbf Z/p\mathbf Z)\xrightarrow{\sim}H^1(S,O_S)^F.
\]
\end{SGAresult}
\textbf{注记 6.11。} 最后一个断言归于 J.-P. Serre [9]。还可对任意 $n$ 的结构群 $\mathbf Z/p^n\mathbf Z$ 发展类似理论，只须以 Witt \Term{group-scheme}{群概形} $W_n$ 代替 $\mathbf G_a$，参见上述文献。注意，在特征 $p>0$ 时，Kummer 理论不再给出关于 $p$ 阶\Term{principal-covering}{主覆盖}的信息，因为此时 $\mu_p$ 为“无穷小”群，即在基上为\Term{radicial}{根性}的，因而与 $\mathbf Z/p\mathbf Z$ 没有直接关系。所以，乍看之下，这些\Term{covering}{覆盖}的理论似乎不能像 6.6 那样由 \Term{picard-scheme}{Picard 概形}的理论处理；为明确起见，可设 $S$ 为\Term{algebraically-closed-field}{代数闭域}上的\Term{proper}{固有}\Term{scheme}{概形}。不过，回忆 $\operatorname{Pic}_{S/K}$\footnote{关于 $\operatorname{Pic}_{S/K}$ 的定义，见 A. Grothendieck，Bourbaki 讨论班第 232 讲（1962 年 2 月）。} 在原点的 Zariski 切空间可认同于 $H^1(S,O_S)$，就会发现：知道\Term{group-scheme}{群概形} $ {}_p\operatorname{Pic}_{S/k}$，即 $\operatorname{Pic}_{S/k}$ 上乘 $p$ 的核，就既知道 $H^1(S,\mathbf Z/p\mathbf Z)$，也知道 $H^1(S,\mu_p)$。这还决定了 $H^1(S,\alpha_p)$，其中 $\alpha_p$ 为素域上的无穷小\Term{group-scheme}{群概形}，是 $F:\mathbf G_a\to\mathbf G_a$ 的核；它也可描述为一维\Term{trivial-bundle}{平凡} $p$-Lie 代数的限制包络代数的谱。事实上，在这里\Term{exact-sequence}{正合列} 4.5 给出
\[
H^1(S,\alpha_p)\simeq
\operatorname{Ker}\bigl(F:H^1(S,O_S)\longrightarrow H^1(S,O_S)\bigr).
\]
更一般地，若 $\alpha_{p^n}$ 表示 $\mathbf G_a$ 中 $F$ 的 $n$ 次迭代的核，则有
\[
H^1(S,\alpha_{p^n})\simeq
\operatorname{Ker}\bigl(F^n:H^1(S,O_S)\longrightarrow H^1(S,O_S)\bigr).
\]
事实上，知道 $ {}_p\operatorname{Pic}_{S/k}$，等价于知道每个被 $p$ 零化的\Term{finite}{有限}交换代数群 $G$ 的 $H^1(S,G)$；更一般地，知道 $ {}_{p^n}\operatorname{Pic}_{S/k}$，等价于知道每个被 $p^n$ 零化的\Term{finite}{有限}交换代数群 $G$ 的 $H^1(S,G)$。这是下述定理的结果；在所考虑的情形下，该定理同时包含 Kummer 理论和 Artin–Schreier 理论。

设 $G$ 为 $k$ 上的\Term{finite}{有限}代数群，
\[
D(G)=\operatorname{Hom}_{k\text{-群}}(G,\mathbf G_m)
\]
为其\Term{cartier-dual}{Cartier 对偶}；其\Term{affine}{仿射}代数的底层向量空间，是 $G$ 的\Term{affine}{仿射}代数的对偶空间，
\SourcePage{254}
即 Dieudonné–Cartier 意义下 $G$ 的超代数。则有典范\Term{isomorphism}{同构}
\[
H^1(S,G)\simeq
\operatorname{Hom}_{k\text{-群}}\bigl(D(G),\operatorname{Pic}_{S/k}\bigr).
\tag{*}
\]
（注意：$S$ 为\Term{algebraically-closed-field}{代数闭域} $k$ 上的\Term{proper}{固有}\Term{scheme}{概形}，且 $H^0(S,O_S)=k$。）这个公式还可以表述为：第 2 节提到的 $S$ 的“真正\Term{fundamental-group}{基本群}”的阿贝尔化，\Term{isomorphism}{同构}于各 $D(P_i)$ 的\Term{projective-limit}{投极限}，其中 $P_i$ 遍历 $\operatorname{Pic}_{S/k}$ 的\Term{finite}{有限}代数子群；将这一\Term{projective-limit}{投极限}记为 $T^\bullet(\operatorname{Pic}_{X/k})$。当 $S$ 为\Term{abelian-variety}{阿贝尔簇}时，由 2.1 已知，该群也\Term{isomorphism}{同构}于“真正的”\Term{tate-module}{Tate 模}
\[
T_{\boldsymbol\cdot}(S)=\varprojlim_n{}_nS.
\]
于是\Term{isomorphism}{同构}（*）可写成更醒目的形式
\[
\operatorname{Ext}^1(A,G)\simeq\operatorname{Hom}(D(G),B),
\]
其中 $A$ 为\Term{abelian-variety}{阿贝尔簇}，$B$ 为其对偶簇，$G$ 为 $k$ 上的\Term{finite}{有限}代数群。刚才指出的这些结果，还可推广到以任意基\Term{prescheme}{预概形}代替 $k$ 的情形，并推广到\Term{finite-group}{有限群}之外的其他系数群 $G$。

\SGAsection{7}{参考文献}
\begin{enumerate}[label={[\arabic*]},leftmargin=2em]
\item A. Grothendieck，关于同调代数的若干问题（\textit{Sur quelques points d'algèbre homologique}），Tôhoku Math. J. 9（1957），119—221 页。
\item A. Grothendieck，带\Term{structure-sheaf}{结构层}的\Term{fibred-space}{纤维空间}的一般理论（\textit{A general theory of fibre spaces with structure sheaf}），University of Kansas，1955。
\item A. Grothendieck，同调扭性与有理\Term{section}{截面}（\textit{Torsion homologique et sections rationnelles}），Chevalley 讨论班，1958 年 6 月 16 日。
\item S. Lang，\Term{abelian-variety}{阿贝尔簇}（\textit{Abelian varieties}），Interscience Tracts in Pure and Applied Mathematics，第 7 号，纽约。
\item J.-P. Serre，代数群与类域（\textit{Groupes algébriques et corps de classes}），Actualités Scientifiques et Industrielles，第 1264 号，Hermann，巴黎，1959。
\item J.-P. Serre，投代数群（\textit{Groupes proalgébriques}），Publications Mathématiques de l'IHES 7（1960），1—67 页。
\item J.-P. Serre，代数\Term{fibred-space}{纤维空间}（\textit{Espaces fibrés algébriques}），Chevalley 讨论班，1958 年 4 月 21 日。
\item J.-P. Serre，特征 $p$ 的\Term{abelian-variety}{阿贝尔簇}的若干性质（\textit{Quelques propriétés des variétés abéliennes en caractéristique p}），Amer. J. Math. 80（1958），715—739 页。
\item J.-P. Serre，关于特征 $p$ 的代数簇的拓扑（\textit{Sur la topologie des variétés algébriques en caractéristique p}），Symposium Internacional de Topologia Algebraica，1958。
\item J.-P. Serre，关于\Term{unirational-variety}{单有理簇}的\Term{fundamental-group}{基本群}（\textit{On the fundamental group of a unirational variety}），J. London Math. Soc. 34（1959），481—484 页。
\end{enumerate}
